\documentclass[12pt]{article}
\usepackage[margin=1in]{geometry}
\usepackage{setspace}
\usepackage{parskip}

\doublespacing

\begin{document}

\noindent \textbf{To:} Ambassador Alexander Kmentt, Director, Disarmament, Arms Control and Non-Proliferation, Austrian Federal Ministry for European and International Affairs \\
\textbf{From:} [Group Members] \\
\textbf{Re:} The Rhetoric-Action Gap in Multilateral Arms Control

\section*{The Problem}

The multilateral arms control regime rests on an implicit ethical bargain: states commit publicly to humanitarian norms around weapons, and in exchange gain legitimacy and access to cooperative security frameworks. We argue that this bargain is breaking down. States increasingly say one thing at the General Assembly podium and do another in the voting booth or on the export market. This memo identifies three dimensions of that gap and proposes policy responses grounded in transparency and accountability.

Our argument draws on an original dataset we constructed by combining the UN General Debate Corpus (9,046 speeches, 1970--2023) with UNGA voting records, V-Dem regime scores, SIPRI arms transfer data, and the texts of major arms control treaties. We use natural language processing to measure how closely each country's rhetoric tracks the language of specific treaties and how those patterns vary by regime type, nuclear status, and arms trade position.

\section*{Initial Findings}

Three patterns stand out from our preliminary analysis:

First, nuclear weapon states and non-nuclear weapon states are talking past each other. The rhetorical distance between these two groups has grown substantially since 2013, when the humanitarian initiative conferences in Oslo and Nayarit began reshaping disarmament discourse. On TPNW-related resolutions, non-nuclear states vote in favor at roughly 90\% while nuclear weapon states vote at 0\%. This is not just a policy disagreement---it reflects a widening gap in how these states frame the ethics of nuclear weapons.

Second, there is a measurable paradox among Non-Aligned Movement autocracies. These states vote pro-disarmament at very high rates but their General Assembly speeches emphasize sovereignty, national liberation, and anti-colonial themes rather than the humanitarian language that characterizes the treaties they claim to support. The gap between their rhetoric and their votes is the largest of any group we examined.

Third, major arms exporters---regardless of whether they are democracies or autocracies---use strikingly similar language when discussing arms control treaties. The United States, Russia, France, China, and the United Kingdom converge rhetorically in ways that cut across their geopolitical rivalries. This suggests material interest in the arms trade shapes diplomatic language more than ideology or regime type does.

\section*{Policy Recommendations}

We plan to develop three recommendations. The first is a transparency mechanism at NPT Review Conferences that would require delegations to reconcile their stated positions with their voting records and treaty participation. The second links Arms Trade Treaty compliance review to consistency between export rhetoric and actual transfer patterns. The third proposes a civil-society-led accountability index---modeled on Transparency International's Corruption Perceptions Index---that would score and publish the rhetoric-action gap for every UN member state annually.

Each of these operates through transparency rather than coercion, which we believe makes them more politically feasible and harder to object to on sovereignty grounds.

\section*{Likely Objections}

The strongest objection is that quantitative text analysis oversimplifies diplomatic language and cannot capture strategic ambiguity. We plan to address this by validating our findings across multiple methods and by acknowledging the limits of computational approaches. A second objection is that naming and shaming does not constrain great powers. We will argue, drawing on the TPNW case, that normative pressure operates on longer timescales and constrains the menu of legitimate options even when it does not produce immediate compliance. A third criticism is that this framework unfairly targets the Global South. Our data shows the gap is not regime-specific---democratic arms exporters are among the worst offenders---and we will design the accountability index to apply symmetrically.

\end{document}
